\chapter{Обзор литературы}
\label{chap:lr}

Многофазная фильтрация в трещиновато-пористых пластах описывается связанной системой нелинейных уравнений для трёх полей: давления, водонасыщенности трещин и матричных блоков. Её решения содержат подвижные резкие фронты вытеснения, которые грубые разностные схемы аппроксимируют неудовлетворительно, а уменьшение шага сетки выводит полномасштабную модель в область неприемлемых вычислительных затрат.

Постановка восстановления детальных полей по грубым приближениям совпадает с задачей сверхразрешения изображений. Сверточные и остаточные архитектуры точны по среднеквадратичной метрике, но сглаживают высокочастотное содержание~\autocite{lim2017enhanced}; состязательное обучение (GAN) восстанавливает резкие границы~\autocite{cvf_srgan}. В нефтяной инженерии суперразрешение уже применялось к полям насыщенности и течения, в том числе с физически обоснованной регуляризацией~\autocite{li2025spe}.

Ключевое ограничение опубликованных работ — применимость только к однопоровым моделям. Задача одновременного восстановления трёх связанных полей для систем двойной пористости в литературе до настоящего времени не рассматривалась, что определяет научную новизну работы.